%\title{University of Bristol Thesis Template}
\RequirePackage[l2tabu]{nag}		% Warns for incorrect (obsolete) LaTeX usage
%
% File: memoirthesis.tex
% Author: Victor Brena
% Description: Contains the thesis template using memoir class,
% which is mainly based on book class but permits better control of 
% chapter styles for example. This template is an adaptation and 
% modification of Oscar's.
% 
% Memoir is a flexible class for typesetting poetry, fiction, 
% non-fiction and mathematical works as books, reports, articles or
% manuscripts. CTAN repository is found at:
% http://www.ctan.org/tex-archive/macros/latex/contrib/memoir/
%
% UoB guidelines for thesis presentation were found at:
% http://www.bris.ac.uk/esu/pg/pgrcop11-12topic.pdf#page=49
%
\documentclass[a4paper,11pt,leqno,oneside]{memoir}
%
\usepackage{colortbl}      % \rowcolor
\usepackage{amsmath}
\usepackage{amssymb}       % \mathbb{I}
\usepackage{amsthm}
\usepackage{algorithm}
\usepackage{algpseudocode}
% Adding metadata:
\usepackage[utf8]{inputenc}
\usepackage[T1]{fontenc}
\usepackage[usenames,dvipsnames,svgnames,table]{xcolor}
\usepackage{calligra}
\usepackage{eso-pic}
\usepackage{tikz}
\usetikzlibrary{arrows,shapes,snakes,
		       automata,backgrounds,
		       petri,topaths}

\usepackage{tcolorbox}

\newtcolorbox{bidentidad}[2][]{
  enhanced,
  skin=enhancedlast jigsaw,
  attach boxed title to top left={xshift=-4mm,yshift=-0.5mm},
  fonttitle=\bfseries\sffamily,
  colbacktitle=darkgray,
  colframe=darkgray,
  interior style={
    top color=blue!3,
    bottom color=red!3
  },
  boxed title style={
    empty,
    arc=0pt,
    outer arc=0pt,
    boxrule=0pt
  },
  underlay boxed title={
    \fill[darkgray] 
      (title.north west) -- 
      (title.north east) -- 
      +(\tcboxedtitleheight-1mm,-\tcboxedtitleheight+1mm) -- 
      ([xshift=4mm,yshift=0.5mm]frame.north east) -- 
      +(0mm,-1mm) -- 
      (title.south west) -- cycle;
    \fill[darkgray] 
      ([yshift=-0.5mm]frame.north west) -- 
      +(-0.4,0) -- 
      +(0,-0.3) -- cycle;
    \fill[darkgray] 
      ([yshift=-0.5mm]frame.north east) -- 
      +(0,-0.3) -- 
      +(0.4,0) -- cycle; 
  },
title = Questions
}


\newtcolorbox{bidentidad2}[2][]{
  enhanced,
  skin=enhancedlast jigsaw,
  attach boxed title to top left={xshift=-4mm,yshift=-0.5mm},
  fonttitle=\bfseries\sffamily,
  colbacktitle=darkgray,
  colframe=darkgray,
  interior style={
    top color=blue!3,
    bottom color=red!3
  },
  boxed title style={
    empty,
    arc=0pt,
    outer arc=0pt,
    boxrule=0pt
  },
  underlay boxed title={
    \fill[darkgray] 
      (title.north west) -- 
      (title.north east) -- 
      +(\tcboxedtitleheight-1mm,-\tcboxedtitleheight+1mm) -- 
      ([xshift=4mm,yshift=0.5mm]frame.north east) -- 
      +(0mm,-1mm) -- 
      (title.south west) -- cycle;
    \fill[darkgray] 
      ([yshift=-0.5mm]frame.north west) -- 
      +(-0.4,0) -- 
      +(0,-0.3) -- cycle;
    \fill[darkgray] 
      ([yshift=-0.5mm]frame.north east) -- 
      +(0,-0.3) -- 
      +(0.4,0) -- cycle; 
  },
title = More generally;
}

\usepackage{enumitem}
\usepackage{bbding}
\usepackage{graphicx}
\usepackage{xcolor}
\usepackage{enumitem}
\usepackage{frcursive}
\usepackage{pifont}
\setlist[itemize]{label=\textcolor{darkgray}{\HandCuffRight}}
\tcbuselibrary{skins}
\usepackage{psvectorian}
\renewcommand*{\psvectorianDefaultColor}{Blue}%
\usepackage{datetime}
\usepackage{ifpdf}
\ifpdf
\pdfinfo{
   /Author (Peterson Yuhala)
   /Title (MII ENSPY)
   /Keywords (Xen; )
   /CreationDate (D:\pdfdate)
}
\fi
\ifdraftdoc 
	\usepackage{draftwatermark}
	\SetWatermarkScale{0.3}
	\SetWatermarkText{\bf Draft: \today}
\fi
%
\newsubfloat{figure}
\newsubfloat{table}
\usepackage{booktabs}
\usepackage{siunitx}

\usetikzlibrary{chains}

\usepackage{multirow}
\usepackage{algorithm}
\usepackage[noend]{algpseudocode}
\makeatletter
\def\BState{\State\hskip-\ALG@thistlm}
\makeatother

% Better page layout for A4 paper
\settrimmedsize{297mm}{210mm}{*}
\setlength{\trimtop}{0pt} 
\setlength{\trimedge}{\stockwidth} 
\addtolength{\trimedge}{-\paperwidth} 
\settypeblocksize{634pt}{448.13pt}{*} 
\setulmargins{4cm}{*}{*} 
\setlrmargins{*}{*}{1.5} 
\setmarginnotes{17pt}{51pt}{\onelineskip} 
\setheadfoot{\onelineskip}{2\onelineskip} 
\setheaderspaces{*}{2\onelineskip}{*} 
\checkandfixthelayout
%
\frenchspacing
\usepackage{fouriernc}
%%% MINITOC SUPPRIMÉ — remplacé par etoc (voir plus bas)
\usepackage{fancybox}
\usepackage{makeidx}


\definecolor{marron}{RGB}{60,30,10}
\definecolor{darkblue}{RGB}{0,0,80}
\definecolor{lightblue}{RGB}{80,80,80}
\definecolor{darkgreen}{RGB}{0,80,0}
\definecolor{darkgray}{RGB}{128,0,0}
\definecolor{darkred}{RGB}{80,0,0}
\definecolor{shadecolor}{rgb}{0.97,0.97,0.97}
\usepackage{fourier-orns}
\usepackage[T1]{fontenc}


\makeatletter
\def\headrule{{\color{darkgray}\raisebox{-2.1pt}[10pt][10pt]{\leafright} \hrulefill \raisebox{-2.1pt}[10pt][10pt]{~~~\decofourleft \decotwo\decofourright~~~} \hrulefill \raisebox{-2.1pt}[10pt][10pt]{ \leafleft}}}
\makeatother
\usepackage{pdflscape} 
\usepackage[
final,
stretch=10,
protrusion=true,
tracking=true,
spacing=on,
kerning=on,
expansion=true]{microtype}

\newcommand{\ornamento}{\vspace{-2em} \noindent \textcolor{darkgray}{\hrulefill~ \raisebox{-2.5pt}[10pt][10pt]{\leafright \decofourleft \decothreeleft  \aldineright \decotwo \floweroneleft \decoone   \floweroneright \decotwo \aldineleft\decothreeright \decofourright \leafleft} ~  \hrulefill}}
\newcommand{\ornpar}{\noindent \textcolor{darkgray}{ \raisebox{-1.9pt}[10pt][10pt]{\leafright} \hrulefill \raisebox{-1.9pt}[10pt][10pt]{\leafright \decofourleft \decothreeleft  \aldineright \decotwo \floweroneleft \decoone}}}
\newcommand{\ornimpar}{\textcolor{darkgray}{\raisebox{-1.9pt}[10pt][10pt]{\decoone \floweroneright \decotwo \aldineleft \decothreeright \decofourright \leafleft} \hrulefill \raisebox{-1.9pt}[10pt][10pt]{\leafleft}}}
%
\OnehalfSpacing 
%
\setsecnumdepth{subsection} 
\maxsecnumdepth{subsubsection}
%
\usepackage{calc,soul,fourier}
\makeatletter 
\newlength\dlf@normtxtw 
\setlength\dlf@normtxtw{\textwidth} 
\newsavebox{\feline@chapter} 
\newcommand\feline@chapter@marker[1][4cm]{%
	\sbox\feline@chapter{% 
		\resizebox{!}{#1}{\fboxsep=1pt%
			\colorbox{darkgray}{\color{white}\thechapter}% 
		}}%
		\rotatebox{90}{% 
			\resizebox{%
				\heightof{\usebox{\feline@chapter}}+\depthof{\usebox{\feline@chapter}}}% 
			{!}{\scshape\so\@chapapp}}\quad%
		\raisebox{\depthof{\usebox{\feline@chapter}}}{\usebox{\feline@chapter}}%
} 
\newcommand\feline@chm[1][4cm]{%
	\sbox\feline@chapter{\feline@chapter@marker[#1]}% 
	\makebox[0pt][c]{% aka \rlap
		\makebox[1cm][r]{\usebox\feline@chapter}%
	}}
\makechapterstyle{daleifmodif}{
	\renewcommand\chapnamefont{\normalfont\Large\scshape\raggedleft\so} 
	\renewcommand\chaptitlefont{\normalfont\Large\bfseries\scshape} 
	\renewcommand\chapternamenum{} \renewcommand\printchaptername{} 
	\renewcommand\printchapternum{\null\hfill\feline@chm[2.5cm]\par} 
	\renewcommand\afterchapternum{\par\vskip\midchapskip} 
	\renewcommand\printchaptertitle[1]{\color{darkgray}\chaptitlefont\raggedleft ##1\par}
} 
\makeatother 

\chapterstyle{daleifmodif}
%
\renewcommand{\sectionmark}[1]{\markright{\thesection~- ~#1}}
\renewcommand{\chaptermark}[1]{\markboth{\chaptername~\thechapter~-~ #1}{}}

\makepagestyle{myvf} 
\makeoddfoot{myvf}{}{\ornimpar \\ \large \hfill \sffamily\bf \textcolor{black}{\leafNE ~~~ Page \thepage}
}{} 
\makeevenfoot{myvf}{}{\ornpar   \\ \large  \sffamily\bf \textcolor{black}{\thepage ~~~ \reflectbox{\leafNE}}  \hfill}{} 
\makeevenhead{myvf}{\leftmark}{\ornamento}{\rightmark} 
\makeoddhead{myvf}{\leftmark}{\ornamento}{\rightmark}
\pagestyle{myvf}

\newcommand{\clearemptydoublepage}{\newpage{\thispagestyle{empty}\cleardoublepage}}
%
\makeindex

\usepackage{import}

\usepackage{lipsum}
\usepackage{subcaption}
\usepackage{amsfonts}
\usepackage[centertags]{amsmath}
\usepackage{stmaryrd}
\usepackage{pgfplots}
\pgfplotsset{yticklabel style={text width=3em,align=right}}

%%% ETOC — remplace minitoc pour les sommaires de chapitres
\usepackage{etoc}

\newcommand{\chaptertoc}{%
  \begingroup
  \etocsettocstyle{\noindent\textbf{Contents}\vskip 4pt\hrule\vskip 4pt}{\vskip 4pt\hrule\bigskip}%
  \etocsetnexttocdepth{2}%
  \localtableofcontents
  \endgroup
}
\usepackage{xltabular}
\usepackage{amssymb}
\usepackage{amsthm}
\usepackage{newlfont}
\usepackage{layouts}
\usepackage{graphicx}
\usepackage{longtable,rotating}

\usepackage{url}
\usepackage{etex}
\usepackage{fixltx2e}
\usepackage[french]{babel}
\usepackage{memhfixc}
\usepackage{color}

\usepackage{enumerate}
\usepackage{footnote}
\usepackage{afterpage}
\usepackage{listings}
\usepackage[object=vectorian]{pgfornament}
\usetikzlibrary{shapes.geometric,calc}
\definecolor{fondpaille}{cmyk}{0,0,0.1,0}
\usepackage{pagecolor}
\usepackage{microtype}
\usepackage{rotfloat}
\usepackage{alltt}

\tcbset{
    Baystyle/.style={
        sharp corners,
        enhanced,
        boxrule=6pt,
        colframe=OliveGreen,
        height=600pt,
        width=500pt,
        borderline={8pt}{-11pt}{},
    }
}

\widowpenalty=1000
\clubpenalty=1000
%
\newcommand{\keywords}[1]{\par\noindent{\small{\bf Keywords:} #1}}
\newcommand{\parcial}[2]{\frac{\partial#1}{\partial#2}}
\newcommand{\vectorr}[1]{\mathbf{#1}}
\newcommand{\vecol}[2]{\left(
	\begin{array}{c} 
		\displaystyle#1 \\
		\displaystyle#2
	\end{array}\right)}
\newcommand{\mados}[4]{\left(
	\begin{array}{cc}
		\displaystyle#1 &\displaystyle #2 \\
		\displaystyle#3 & \displaystyle#4
	\end{array}\right)}
\newcommand{\pgftextcircled}[1]{
    \setbox0=\hbox{#1}%
    \dimen0\wd0%
    \divide\dimen0 by 2%
    \begin{tikzpicture}[baseline=(a.base)]%
        \useasboundingbox (-\the\dimen0,0pt) rectangle (\the\dimen0,1pt);
        \node[circle,draw,outer sep=0pt,inner sep=0.1ex] (a) {#1};
    \end{tikzpicture}
}
\newcommand{\ra}[1]{\renewcommand{\arraystretch}{#1}}
\newcommand{\myrowcolour}{\rowcolor[gray]{0.925}}

\newcommand{\range}[1]{\textnormal{range }#1}
\newcommand{\innerp}[2]{\left\langle#1,#2\right\rangle}
\newcommand{\prom}[1]{\left\langle#1\right\rangle}
\newcommand{\tra}[1]{\textnormal{tra} \: #1}
\newcommand{\sign}[1]{\textnormal{sign\,}#1}
\newcommand{\sech}[1]{\textnormal{sech} #1}
\newcommand{\diag}[1]{\textnormal{diag} #1}
\newcommand{\arcsech}[1]{\textnormal{arcsech} #1}
\newcommand{\arctanh}[1]{\textnormal{arctanh} #1}

\newcommand{\blackged}{\hfill$\blacksquare$}
\newcommand{\whiteged}{\hfill$\square$}
\newcounter{proofcount}
\renewenvironment{proof}[1][\proofname.]{\par
 \ifnum \theproofcount>0 \pushQED{\whiteged} \else \pushQED{\blackged} \fi%
 \refstepcounter{proofcount}
 \normalfont 
 \trivlist
 \item[\hskip\labelsep
       \itshape
   {\bf\em #1}]\ignorespaces
}{%
 \addtocounter{proofcount}{-1}
 \popQED\endtrivlist
}
%
\let\oldsqrt\sqrt
\def\sqrt{\mathpalette\DHLhksqrt}
\def\DHLhksqrt#1#2{%
\setbox0=\hbox{$#1\oldsqrt{#2\,}$}\dimen0=\ht0
\advance\dimen0-0.2\ht0
\setbox2=\hbox{\vrule height\ht0 depth -\dimen0}%
{\box0\lower0.4pt\box2}}
%
\newcommand{\mycaption}[2][\@empty]{
	\captionnamefont{\scshape} 
	\changecaptionwidth
	\captionwidth{0.9\linewidth}
	\captiondelim{.\:} 
	\indentcaption{0.75cm}
	\captionstyle[\centering]{}
	\setlength{\belowcaptionskip}{10pt}
	\ifx \@empty#1 \caption{#2}\else \caption[#1]{#2}
}
\newcommand*\circled[1]{\tikz[baseline]{\node[shape=circle,draw,inner sep=2pt] (char) {#1};}}

\newcommand{\mysubcaption}[2][\@empty]{
	\subcaptionsize{\small}
	\hangsubcaption
	\subcaptionlabelfont{\rmfamily}
	\sidecapstyle{\raggedright}
	\setlength{\belowcaptionskip}{10pt}
	\ifx \@empty#1 \subcaption{#2}\else \subcaption[#1]{#2}
}
%
\newcommand{\sectionlinetwo}[2]{%
  \nointerlineskip \vspace{.5\baselineskip}\hspace{\fill}
  {\color{#1}
    \resizebox{0.5\linewidth}{2ex}
    {{%
    {\begin{tikzpicture}
    \node  (C) at (0,0) {};
    \node (D) at (9,0) {};
    \path (C) to [ornament=#2] (D);
    \end{tikzpicture}}}}}%
    \hspace{\fill}
    \par\nointerlineskip \vspace{.5\baselineskip}
  }
%
\usepackage{Carrickc,ArtNouvc,Kramer,Konanur,Rothdn}
\usepackage{lettrine}
\newcommand{\initial}[1]{%
	\lettrine[lines=3,lhang=0.33,nindent=0em]{
		\color{darkgray}
     		{\textsc{#1}}}{}}
     		
\newcommand{\corner}[1]{%
  \begin{tikzpicture}[color=darkgray,remember picture, overlay]
    \node[anchor=north west] at (current page.north west){%
      \pgfornament[width=2cm]{#1}};
    \node[anchor=north east] at (current page.north east){%
      \pgfornament[width=2cm,symmetry=v]{#1}};
    \node[anchor=south west] at (current page.south west){%
      \pgfornament[width=2cm,symmetry=h]{#1}};
    \node[anchor=south east] at (current page.south east){%
      \pgfornament[width=2cm,symmetry=c]{#1}};
  \end{tikzpicture}%
}
\newcommand{\cornerplus}[2]{%
  \begin{tikzpicture}[color=darkgray,remember picture, overlay]
    \node[anchor=north west] at (current page.north west){%
      \pgfornament[width=2cm]{#1}};
    \node[anchor=north east] at (current page.north east){%
      \pgfornament[width=2cm,symmetry=v]{#1}};
    \node[anchor=south west] at (current page.south west){%
      \pgfornament[width=2cm,symmetry=h]{#1}};
    \node[anchor=south east] at (current page.south east){%
      \pgfornament[width=2cm,symmetry=c]{#1}};
    \node[anchor=north] at (current page.north){%
      \pgfornament[width=6.5cm,symmetry=h]{#2}};
    \node[anchor=south] at (current page.south){%
      \pgfornament[width=6.5cm]{#2}};
  \end{tikzpicture}%
}
\newcommand{\pt}[1]{%
  \begin{tikzpicture}[color=darkgray,remember picture, overlay]
    \node[anchor=north] at (current page.north){%
      \pgfornament[symmetry=h]{#1}};
    \node[anchor=south] at (current page.south){%
      \pgfornament[symmetry=h]{#1}};
  \end{tikzpicture}%
}
%
\AtBeginDocument{\renewcommand{\bibname}{R\'ef\'erences bibliographiques}}
\AtBeginDocument{\renewcommand{\refname}{R\'ef\'erences bibliographiques}}

\theoremstyle{plain}
\newtheorem{theo}{Theorem}[chapter]
\theoremstyle{plain}
\newtheorem{prop}{Proposition}[chapter]
\theoremstyle{plain}
\theoremstyle{definition}
\newtheorem{dfn}{Definition}[chapter]
\theoremstyle{plain}
\newtheorem{lema}{Lemma}[chapter]
\theoremstyle{plain}
\newtheorem{cor}{Corollary}[chapter]
\theoremstyle{plain}
\newtheorem{resu}{Result}[chapter]
%
\hyphenation{res-pec-tively}
\hyphenation{mono-ti-ca-lly}
\hyphenation{hypo-the-sis}
\hyphenation{para-me-ters}
\hyphenation{sol-va-bi-li-ty}
%
\usepackage[acronym]{glossaries}
\makeglossary

\newacronym{enspy}{ENSPY}{\'{E}cole Nationale Sup\'{e}rieure Polytechnique Yaounde}
\newacronym{uy1}{UY1}{Université de Yaounde I}
\newacronym{imt}{IMT}{IMT Atlantique}
\newacronym{ls2n}{LS2N}{Laboratoire des Sciences du Numérique de Nantes}
\newacronym{iota}{IOTA}{IOT \& Applications}

\newglossaryentry{bench}
{
	name = Benchmark,
    description = {In computer science, a benchmark is a test to measure the performance of a system to compare it to others},
    plural = benchmark
}

\newglossaryentry{bus}
{
	name = Bus,
     description = {In computer science, a bus is a communication system that transfers data between components inside a computer, or between computers.},
    plural = bus
}

\newglossaryentry{ops}
{
    name=Open source,
    description={The open source designation, or "open source code", applies to software whose license meets criteria precisely established by the Open Source Initiative, that is, the possibilities of free redistribution, access to source code and the creation of derivative works.},
    plural=open source
}

\newglossaryentry{interrupt}
{
	name = Interrupt,
    description = {An interrupt is a signal from a device attached to a computer or from a program within the computer that requires the operating system to stop and figure out what to do next},
    plural = interrupt
}

\newglossaryentry{sc}
{
	name = Source code,
    description = {A source code is a text that represents a program instructions as written by a programmer.},
    plural = source code
}


\usepackage{chngcntr}
\usepackage[section]{placeins}  
\usepackage{titletoc}

% ===================== HYPERREF (chargé en dernier) =====================
\usepackage[colorlinks=true,
            linkcolor=black,
            citecolor=black,
            urlcolor=blue,
            linktoc=all,
            pdfusetitle]{hyperref}
% =======================================================================

\begin{document}
\sloppy

\lstloadlanguages{bash, Java, sql, python}
\lstset{basicstyle=\small,
    keywordstyle=\color{blue}\ttfamily,
    identifierstyle=\ttfamily,
    commentstyle=\color{green}\itshape\ttfamily,
    stringstyle=\sffamily,
    showstringspaces=false,
    numbers=left,
    numberstyle=\tiny}

\frontmatter
\pagenumbering{roman}
%
%
% File: Title.tex
% Author: V?ctor Bre?a-Medina
% Description: Contains the title page
%
% UoB guidelines:
% 
% At the top of the title page, within the margins, the dissertation should give the title and, if 
% necessary, sub-title and volume number. If the dissertation is in a language other than English, the 
% title must be given in that language and in English. The full name of the author should be in the 
% centre of the page. At the bottom centre should be the words ?A dissertation submitted to the 
% University of Bristol in accordance with the requirements for award of the degree of ? in the 
% Faculty of ...?, with the name of the school and month and year of submission. The word count of 
% the dissertation (text only) should be entered at the bottom right-hand side of the page.
%
%
%\hspace{5cm}
%\pagecolor{fondpaille}%\afterpage{\nopagecolor}

\begin{titlingpage}
\begin{SingleSpace}

%\calccentering{\unitlength} 
\begin{adjustwidth*}{\unitlength}{-\unitlength}
%Front page corner design...read ornaments documentation for more info
%Peterson Yuhala
%Drawn via eso-pic's shipout hook (fires once, on the very next physical
%page shipped) instead of a plain overlay: this way the frame is tied to
%the actual page 1 output, regardless of where in the flow it's written
%or whether the content below happens to overflow onto a second page.
\AddToShipoutPictureBG*{%
\begin{tikzpicture}[remember picture, overlay, start chain, node distance=-2mm,color=darkgray]
\node (nworn) [shift={(5mm,-5mm)}, anchor=north west, on chain ] at (current page.north west) {\pgfornament[width=10mm]{19}};
\foreach \i in {1,...,17}
\node [on chain] {\pgfornament[width=10mm]{19}};
\node (neorn) [on chain] {\pgfornament[width=10mm]{19}};
\foreach \i in {1,...,25}
\node [continue chain=going below, on chain] {\pgfornament[width=10mm]{19}};
\node (seorn) [on chain] {\pgfornament[width=10mm]{19}};
\foreach \i in {1,...,17}
\node [continue chain=going left, on chain] {\pgfornament[width=10mm]{19}};
\node (sworn) [on chain] {\pgfornament[width=10mm]{19}};
\foreach \i in {1,...,25}
\node [continue chain=going above, on chain] {\pgfornament[width=10mm]{19}};
\end{tikzpicture}%
}
\vspace*{13mm}
\vspace{-4cm}
\begin{center}
 	\begin{tabular}{ccc}
				{\large UNIVERSITE DE YAOUNDE I} & \multirow{8}{*}{\includegraphics[ width =3cm , height =3cm]{logos/uy1.png}} & \large{UNIVERSITY OF YAOUNDE I}\\
				%& & \\
				\textbf{*******} &    & \textbf{*******} \\
				%& & \\
				\textcolor{darkgray}{ECOLE NATIONALE SUPERIEURE}& & \textcolor{darkgray}{NATIONAL ADVANCED SCHOOL}\\ 
				\textcolor{darkgray}{POLYTECHNIQUE DE YAOUNDE}	  & & \textcolor{darkgray}{OF ENGINEERING OF YAOUNDE} \\
				%& & \\
				\textbf{********}    & & \textbf{*******} \\
				%& & \\
				DEPARTEMENT DE GENIE& & DEPARTMENT OF COMPUTER \\ 
				INFORMATIQUE & & ENGINEERING \\
				\end{tabular}
	
\end{center}
\hspace{3cm}
\vspace{2.3cm}
\begin{center}
\rule[0.5ex]{\linewidth}{2pt}\vspace*{-\baselineskip}\vspace*{3.2pt}
\rule[0.5ex]{\linewidth}{1pt}\\[\baselineskip]
{\LARGE \textcolor{darkgray}{\textsc{Modélisation d'un module d'extraction automatique de données pour documents administratifs : étude de cas Dolibarr}}}\\[4mm]

% {\Large \textcolor{darkgray}{Applied to the domain 0 in the Xen virtualization system}}\\
\rule[0.5ex]{\linewidth}{1pt}\vspace*{-\baselineskip}\vspace{3.2pt}
\rule[0.5ex]{\linewidth}{2pt}\\
\vspace{10mm}
{\Large \textbf{Rapport de stage PR\'E-ING\'ENIEUR}}\\
\vspace{15mm}
{\Large Pr\'esent\'e par: } \\
\vspace{4mm}
{\large \textsc{\textbf{\textcolor{darkgray}{FOFACK ALEMDJOU HENRI JOEL}}}}\\
\vspace{15mm}
{\Large Matricule: } \\
\vspace{4mm}
{\large \textsc{\textbf{\textcolor{darkgray}{22P231}}}}\\
\vspace{3.5mm}
\vspace{20mm}
{\Large Sous la supervision de :}\\
\vspace{4mm}
{\normalsize \textsc{\textbf{\textcolor{darkgray}{M.BRELL SANWOUO, PHD Student et CEO de NOVALITIX}}}}\\
\vspace{4.5mm}
\vspace{15mm}


\begin{tabular}{c}
{\Large \textcolor{darkgray}{Ann\'ee acad\'emique 2025-2026}}\\
% {\Large \textcolor{darkgray}{Defended the 08th September 2017}}
\end{tabular}

% \begin{tikzpicture}[remember picture, overlay]
%  \begin{scope}[shift={(current page.south west)},shift={(1,1)},scale=1]
%  \shade[ball color=darkgray,opacity=.6] (0,0) circle (10ex);
%  \shade[ball color=darkgray,opacity=.8] (1.7,1) circle (6ex);
%  \shade[ball color=darkgray,opacity=.8] (1.5,3) circle (2ex);
%  \shade[ball color=darkgray,opacity=.5] (-0.5,3) circle (1ex);
%  \shade[ball color=darkgray,opacity=.8] (1,4) circle (1ex);
%  \shade[ball color=darkgray,opacity=.6] (3.5,2.5) circle (2ex);
%  \shade[ball color=darkgray,opacity=.8] (2.5,3.5) circle (2ex);
%  \end{scope}
%  \end{tikzpicture}
%%%%%%%%%%%%%%%%%%%%%%%%%%%%%%%%%%%%%%%%%%%%%%%%%%%%%%%%%%%%

%\cornerplus{41}{88}
\end{center}


\end{adjustwidth*}

\end{SingleSpace}
\end{titlingpage}

\cleardoublepage

%
% File: Title.tex
% Author: V?ctor Bre?a-Medina
% Description: Contains the title page
%
% UoB guidelines:
% 
% At the top of the title page, within the margins, the dissertation should give the title and, if 
% necessary, sub-title and volume number. If the dissertation is in a language other than English, the 
% title must be given in that language and in English. The full name of the author should be in the 
% centre of the page. At the bottom centre should be the words ?A dissertation submitted to the 
% University of Bristol in accordance with the requirements for award of the degree of ? in the 
% Faculty of ...?, with the name of the school and month and year of submission. The word count of 
% the dissertation (text only) should be entered at the bottom right-hand side of the page.
%
%
%\hspace{5cm}
%\pagecolor{fondpaille}\afterpage{\nopagecolor}
\begin{titlingpage}
\begin{SingleSpace}

%\calccentering{\unitlength} 
\begin{adjustwidth*}{\unitlength}{-\unitlength}
\begin{figure*}
\centering
\includegraphics[height=3cm]{logos/po.png}\hspace{2cm}\includegraphics[height=3cm]{logos/NovalitixLogo.png}
\end{figure*}
\hspace{3cm}
\vspace{0.5cm}
\begin{center}
\rule[0.5ex]{\linewidth}{2pt}\vspace*{-\baselineskip}\vspace*{3.2pt}
\rule[0.5ex]{\linewidth}{1pt}\\[\baselineskip]
{\LARGE \textcolor{darkgray}{\textsc{Modélisation d'un module d'extraction automatique de données pour documents administratifs : étude de cas Dolibarr}}}\\[4mm]
% {\Large \textcolor{darkgray}{Applied to the domain 0 in the Xen virtualization system}}\\
\rule[0.5ex]{\linewidth}{1pt}\vspace*{-\baselineskip}\vspace{3.2pt}
\rule[0.5ex]{\linewidth}{2pt}\\
\vspace{6mm}
{\Large \textbf{Rapport de stage PR\'E-ING\'ENIEUR}}\\
\vspace{15mm}
{\Large Pr\'esent\'e par: } \\
\vspace{4mm}
{\large \textsc{\textbf{\textcolor{darkgray}{FOFACK ALEMDJOU HENRI JOEL}}}}\\
\vspace{15mm}
{\Large Matricule: } \\
\vspace{4mm}
{\large \textsc{\textbf{\textcolor{darkgray}{22P231}}}}\\
\vspace{3.5mm}
\vspace{20mm}
{\Large Sous la supervision de :}\\
\vspace{4mm}
{\normalsize \textsc{\textbf{\textcolor{darkgray}{M.BRELL SANWOUO, PHD Student Et CEO De NOVALITIX}}}}\\
\vspace{4.5mm}
\vspace{15mm}


\begin{tabular}{c}
{\Large \textcolor{darkgray}{Ann\'ee acad\'emique 2025-2026}}\\
% {\Large \textcolor{darkgray}{Defended the 08th September 2017}}
\end{tabular}
\end{center}

\end{adjustwidth*}

\end{SingleSpace}
\end{titlingpage}

\cleardoublepage
%
% Dédicace
\cleardoublepage
\phantomsection
\addcontentsline{toc}{chapter}{Dédicace}
\chapter*{D\'edicace}
%\begin{SingleSpace}
\vspace{3cm}
\begin{center}
  \begin{tikzpicture}[color=darkgray,every node/.style={inner sep=0pt}]
    % Noeud de texte central (taille agrandie, deux lignes)
    \node[text width=12cm,align=center](Text){%
      {\calligra\fontsize{28}{34}\selectfont 
       A  Mes  Tr\`es  Chers  Parents\\
       Et   Fr\`eres  Et  Soeurs}
    };

    % Coins positionnés à l'extérieur du texte
    \node[anchor=south east](CNW) at (Text.north west)
                   {\pgfornament[width=2cm]{61}};
    \node[anchor=south west](CNE) at (Text.north east)
                   {\pgfornament[width=2cm,symmetry=v]{61}};
    \node[anchor=north east](CSW) at (Text.south west)
                   {\pgfornament[width=2cm,symmetry=h]{61}};
    \node[anchor=north west](CSE) at (Text.south east)
                   {\pgfornament[width=2cm,symmetry=c]{61}};
  \end{tikzpicture}
\end{center}
%\end{SingleSpace}
\clearpage




\clearemptydoublepage
%
% Déclaration
\cleardoublepage
\phantomsection
\addcontentsline{toc}{chapter}{Déclaration}
%
% File: declaration.tex
% Author: V?ctor Bre?a-Medina
% Description: Contains the declaration page
%
% UoB guidelines:
%
% Author's declaration
%
% I declare that the work in this dissertation was carried out in accordance
% with the requirements of the University's Regulations and Code of Practice
% for Research Degree Programmes and that it has not been submitted for any
% other academic award. Except where indicated by specific reference in the
% text, the work is the candidate's own work. Work done in collaboration with,
% or with the assistance of, others, is indicated as such. Any views expressed
% in the dissertation are those of the author.
%
% SIGNED: .............................................................
% DATE:..........................
%
\chapter*{Remerciements}
\begin{SingleSpace}
\begin{quote}
\initial{J}e tiens tout d'abord à remercier le seigneur Dieu tout puissant. Cependant, ce travail n'aurait pas été possible sans l'aide ni l'apport de :
\vspace{6mm}
\begin{itemize}[label=\textcolor{darkgray}{\scalebox{1.2}{\HandCuffRight}}]
	\item \textbf{Pr. Jean-Marc MENAUD}, pour ce stage, l'encadrement et le temps accordé pour nos réflexions;
    \item \textbf{Pr. Bernab\'{e} BATCHAKUI}, pour la rigueur inculquée et la motivation sans cesse renouvelée de travailler;
    \item \textbf{Dr. Anne-Marie CHANA}, pour la discipline appliquée, le suivi et de m'avoir permis participer aux activités du laboratoire IOT \& Applications;
    \item Mes camarades de stage, pour leur amitié, les fous rires et qui constituent aujourd'hui une famille;
    \item Ma famille pour leur appui et leur soutien inébranlable 

    
\end{itemize}

\end{quote}
\end{SingleSpace}
\clearpage

\clearemptydoublepage
%
% Résumé
\cleardoublepage
\phantomsection
\addcontentsline{toc}{chapter}{Résumé}
%
% File: abstract.tex
% Author: V?ctor Bre?a-Medina
% Description: Contains the text for thesis abstract
%
% UoB guidelines:
%
% Each copy must include an abstract or summary of the dissertation in not
% more than 300 words, on one side of A4, which should be single-spaced in a
% font size in the range 10 to 12. If the dissertation is in a language other
% than English, an abstract in that language and an abstract in English must
% be included.

\chapter*{R\'{e}sum\'{e}}
\initial{D}e nos jours, dans les entreprises, nous sommes amen\'{e}s \`{a} g\'{e}rer plusieurs types de documents (contrats, factures, bulletins de paie, CV). L'extraction de donn\'{e}es issues de cette vari\'{e}t\'{e} de documents repr\'{e}sente un d\'{e}fi majeur pour les syst\`{e}mes d'information. En effet, les solutions actuelles manquent d'efficacit\'{e}~: les approches bas\'{e}es sur des r\`{e}gles strictes sont rigides, tandis que les mod\`{e}les d'intelligence artificielle sacrifient souvent la pr\'{e}cision sur les valeurs num\'{e}riques et le contexte positionnel. Pour pallier \`{a} cela, notre recherche met en \oe{}uvre une m\'{e}thodologie hybride innovante qui combine l'extraction par templates, la reconnaissance d'entit\'{e}s nomm\'{e}es (NER) et les techniques de reconnaissance de motifs. Dans un cadre agile, notre syst\`{e}me int\`{e}gre un m\'{e}canisme d'apprentissage continu qui am\'{e}liore significativement la pr\'{e}cision d'extraction au fil des it\'{e}rations m\'{e}tier, en s'adaptant aux nouveaux formats de documents sans n\'{e}cessiter de r\'{e}entra\^{\i}nement lourd. Notre prototype, int\'{e}gr\'{e} comme module \`{a} DOLibaar, a r\'{e}duit la saisie manuelle de 80~\% et atteint plus de 95~\% de pr\'{e}cision d'extraction sur un p\'{e}rim\`{e}tre de 10\,000 documents tests. En comparaison, une approche par expressions r\'{e}guli\`{e}res seules plafonne \`{a} 72~\% de pr\'{e}cision, tandis qu'un mod\`{e}le LLM g\'{e}n\'{e}rique non fine-tun\'{e} atteint \`{a} peine 81~\% sur le m\^{e}me jeu de donn\'{e}es, d\'{e}montrant ainsi la sup\'{e}riorit\'{e} de notre m\'{e}thode hybride. Il s'int\`{e}gre dans les syst\`{e}mes m\'{e}tier via des formats standards (JSON, XML), pour un temps de traitement moyen inf\'{e}rieur \`{a} 500~ms par fichier. Bien que la gestion des documents scann\'{e}s de faible qualit\'{e} reste une limite identifi\'{e}e, notre architecture modulaire ouvre la voie \`{a} l'ajout d'un pr\'{e}traitement OCR, consolidant ainsi l'industrialisation de la solution.

\paragraph{Mot cl\'{e}s : Extraction automatique de donn\'{e}es, Documents administratifs, Reconnaissance d'entit\'{e}s nomm\'{e}es (NER), M\'{e}thodologie hybride, Apprentissage continu, DOLibaar.}


\clearpage


\clearemptydoublepage
%
% Abstract
\cleardoublepage
\phantomsection
\addcontentsline{toc}{chapter}{Abstract}
%
% File: abstract.tex
% Author: V?ctor Bre?a-Medina
% Description: Contains the text for thesis abstract
%
% UoB guidelines:
%
% Each copy must include an abstract or summary of the dissertation in not
% more than 300 words, on one side of A4, which should be single-spaced in a
% font size in the range 10 to 12. If the dissertation is in a language other
% than English, an abstract in that language and an abstract in English must
% be included.

\chapter*{Abstract}
\initial{T}oday's enterprises handle a wide variety of documents, including contracts, invoices, payslips and CVs. Extracting data from such heterogeneous documents is a major challenge for information systems. Existing solutions are often inadequate: rule-based approaches are rigid, while artificial intelligence models frequently sacrifice accuracy on numerical values and positional context. To address these limitations, this research introduces an innovative hybrid methodology that combines template-based extraction, named entity recognition (NER), and pattern recognition techniques. Within an agile framework, the proposed system incorporates a continuous learning mechanism that significantly improves extraction accuracy over successive business iterations, adapting to new document formats without requiring extensive retraining. The prototype, integrated as a module into DOLibaar, reduced manual data entry by 80\% and achieved over 95\% extraction accuracy on a test set of 10,000 documents. In comparison, a regex-only approach plateaus at 72\% accuracy, while a generic, non-fine-tuned LLM reaches only 81\% on the same dataset, demonstrating the superiority of the hybrid method. The solution integrates into business systems through standard formats such as JSON and XML, with an average processing time of less than 500~ms per file. Although handling low-quality scanned documents remains an identified limitation, the modular architecture paves the way for adding OCR preprocessing, thereby supporting the industrialization of the solution.
\paragraph{Keywords: Automatic data extraction, Administrative documents, Named Entity Recognition (NER), Hybrid methodology, Continuous learning, DOLibaar.}

\clearpage

\clearemptydoublepage
%
% Glossaire (sigles et abréviations)
\cleardoublepage
\phantomsection
\addcontentsline{toc}{chapter}{Sigles et Abréviations}
\chapter*{Abréviations Et Sigles }
\begin{table}[H]
    \centering
    \begin{tabular}{p{2cm} p{10cm}}
        \textbf{ENSPY} & École Nationale Supérieure Polytechnique de Yaoundé \\  
        \textbf{UY1} & Université de Yaounde I \\ 
        \textbf{IMT} & Institut Mines Télécoms Atlantique de Nantes\\ 
        \textbf{LS2N} & Laboratoire des Sciences du Numérique de Nantes \\ 
        \textbf{IOTA} & Internet Of Things \& Applications \\ 
        \textbf{API} & Application Programming Interface \\
        \textbf{SSH} & Secure Shell \\
        \textbf{YAML} & Yet Another Markup Language \\
    \end{tabular}
\end{table}

\newpage

\chapter*{Glossaire}
\begin{description}
    \item[API] : (Application Programming Interface) est un programme permettant \`a deux applications distinctes de communiquer et d'\'echanger des donn\'ees.
    \item[Conteneur] : environnement d'ex\'ecution \`a une image Docker.
    \item[Docker] : est un logiciel de virtualisation, plus pr\'ecis\'ement un isolateur permettant de lancer des conteneurs encapsulant des logiciels.
    \item[Dockerfile] : fichier texte contenant des commandes permettant la compilation d'une image Docker.
    \item[DPE] : (Diagnostic de performance \'energ\'etique) donne des informations sur la performance \'energ\'etique et climatique d'un logement ou d'un b\^atiment.
    \item[Endpoint] : est l'extr\'emit\'e d'un canal de communication, point par lequel une API entre en contact avec d'autres syst\`emes.
    \item[ExpressJS] : est un framework pour construire des applications web bas\'e sur NodeJS.
    \item[HTTPS] : (Hypertext Transfer Protocol Secure) est un protocole s\'ecuris\'e d'\'echange de
donn\'ees entre un navigateur web et un site web.
    \item[Image Docker] : r\'esultat que l'on obtient apr\`es la compilation d'un Dockerfile, repr\'esentant notre logiciel.
    \item[Kubernetes] : est un syst\`eme open-source permettant d'automatiser le d\'eploiement, la
gestion de conteneurs d'applications et de clusters.
    \item[NodeJS] : environnement d'ex\'ecution de programme JavaScript en dehors du navigateur
gr\^ace \`a l'encapsulation du moteur V8 de Google Chrome.
    \item[PostgreSQL] : est un syst\`eme de gestion de base de donn\'ees relationnelle et objet.
    \item[SSH] : (Secure Shell) est un protocole s\'ecuris\'e qui impose un \'echange de cl\'es de chiffrement en d\'ebut de connexion entre plusieurs ordinateurs via l'\'echange de paquets authentifi\'es.
    \item[WireGuard] : est un protocole de communication et un logiciel libre et open source permettant de cr\'eer un r\'eseau virtuel.
    \item[YAML] : (Yet Another Markup Language) est un format de repr\'esentation de donn\'ees par
s\'erialisation Unicode.
\end{description}


\clearemptydoublepage

%
\renewcommand{\contentsname}{Table des Mati\`eres}
\maxtocdepth{subsection}

%%% MINITOC SUPPRIMÉ — plus besoin de \dominitoc ni \mtcsetdepth
% Table des matières
\phantomsection
\tableofcontents
\markboth{TABLE DES MATI\`ERES}{}

% Liste des figures
\cleardoublepage
\phantomsection
% Liste des figures
\cleardoublepage
\phantomsection            % Nécessite hyperref, sinon à supprimer
\addtocontents{lof}{\par\nobreak\textbf{\scshape Figure \hfill Page}\par\nobreak}  % Déplacé AVANT
\listoffigures             % Sans l'astérisque
\cleardoublepage

% Liste des tableaux
\cleardoublepage
\phantomsection            % Nécessite hyperref, sinon à supprimer
\addtocontents{lot}{\par\nobreak\textbf{\scshape Tableau \hfill Page}\par\nobreak} % lot ET déplacé AVANT
\listoftables              % Sans l'astérisque
\cleardoublepage           % Remplacé ici
\clearemptydoublepage
%

\clearemptydoublepage
%

\mainmatter
%
\import{chapters/chapter01/}{chap01.tex}
\import{chapters/chapter03/}{chap03.tex}
\import{chapters/chapter04/}{chap04.tex}
%
\backmatter
\bibliographystyle{plain}
\refstepcounter{chapter}
\markboth{R\'EF\'ERENCES BIBLIOGRAPHIQUES}{}
\bibliography{biblio}
\clearemptydoublepage
%
\clearemptydoublepage
%   
\end{document}